% ===================================================================
% 章节模板（精读章）—— 复制本文件到 part/<领域>/<论文别名>.tex 后开写
% 体例说明见 CONTRIBUTING.md；本篇在本领域的坐标先看同目录 partinfo.tex
% ===================================================================
\chapter{论文标题（作者，期刊 年份）}
\label{ch:CHANGEME}
% 主笔：<姓名>；初稿：<日期>
% 原文出处：doi:xxx（所据版本：已刊版 / NBER wXXXX / arXiv vN——数字以所据版本为准）
% partinfo 定位：<一句话，说明本篇在领域脉络中的位置>

\section{这篇论文在说什么}
% 三段以内：当时面对的问题与争论（这条脉络上别人走到了哪里）、论文的回答、
% 为什么值得精读。结尾给论文骨架表（原文章节 | 页码 | 讲什么）。

\section{预备知识}
% 视需要。读者设定：仅具备本科中级微观经济学与计量经济学基础。
% 本章要用而前面章节没讲过的概念（计量概念、制度背景、乃至农学概念），先补齐再用。

\section{数据与测度}
% 数据来源与样本口径；关键变量怎么构造（气候变量、产出变量、适应能力指标）。
% 每个有选择余地的口径都说明论文为什么这么选；易错口径用 pitfall 盒子。
% 气候变量的定义（阈值、度日区间、数据源）逐项写清——跨研究不可比的头号来源。

\section{模型设定与识别策略}
% 逐个方程：先说人话，再上公式，再逐符号解读，再用例子算。
% 识别假设单独写明，并讨论：假设不成立时，系数会被解释成什么。
% 设一个贯穿全章的具体例子（论文的某张主表，或一个假想的县、一种作物），
% 跟到底，不许中途丢掉。
% ident 盒子放可直接搬用的设定；intuition 盒子放经济直觉与设计直觉。

% \begin{ident}[标题]
% \end{ident}
% \begin{pitfall}[标题]
% \end{pitfall}
% \begin{intuition}[标题]
% \end{intuition}

\section{结果与解释}
% 主结果表怎么读（不只抄数字与星号）；系数的经济含义与量级必讲；
% 稳健性检验逐组给结论；跨列/跨表的口径差异要标明。

\section{机制与讨论}
% 视需要。机制检验、异质性（灌溉、品种、地块规模等）、适应渠道；
% 论文的解释与竞争性解释并列。

\section{局限与批评}
% 论文自己承认的局限；后续文献指出的批评；本讲义的评注（必须标明「本讲义的评注」）。

\section{在本领域脉络中的位置}
% 回答了什么旧问题、用了谁的路线、被谁继承、开了什么新问题。
% 指向 partinfo 与同 part 其他章用 \S\ref{} 或 \chref{}。
% 定稿前自查：本 part 的 partinfo 时间线与速查表里，能找到本章了吗？
% 能找到就去 partinfo 把对应「待补」格换成 \chref{ch:CHANGEME}；找不到就先想清楚再交。

% 章末参考文献（只列本章引用的条目），勿删
\chapterrefs
